\section{Hypothesis 1: the sampler}
\label{sec:sampler}

\paragraph{Objective novelty.} On OLMo-2-13B, whose training corpus is
public, we ran a fully paired grid of 3 prompts $\times$ 5 seeds $\times$
\{min-$p$ baseline, $\lambda=1$, $\lambda=2$\} and counted, with infini-gram,
which $n$-grams of each generation occur anywhere in the corpus
(Figure~\ref{fig:phase1}). Novel 4-grams rise from 21.5\% (baseline) to
36.0\% ($\lambda=1$) and 45.5\% ($\lambda=2$): paired bootstrap $\Delta$
+14.6\,pp \ci{7.9}{21.9} and +24.0\,pp \ci{18.2}{29.5}, monotone in $\lambda$
and present in each prompt separately. The probability that a generation
contains a verbatim training block of eight or more words falls from 0.80 to
0.20 in both machine arms, a fourfold reduction stated as a rate (a pilot's
``zero blocks'' was one cell, not the distribution). The coherence cost is
about +1.2 perplexity under a cross-family judge model, uncorrelated with
novelty within arms (Spearman $+0.12$). Three escape modes appeared and were
mapped: recitation (a genre fork into a memorized text), collage, and factual
paraphrase; the entropy \emph{ceiling} of the band, above which the model is
at a genre fork and pushing produces collapse, transfers across model
families. A recitation detector based on the entropy drop between halves of a
generation is reported descriptively only (AUC 0.71 for $\ge$8-word blocks,
precision 0.33 at the published threshold; no false-positive rate could be
measured because the baseline arm lacked telemetry).

\begin{figure}[t]
\centering
\includegraphics[width=0.72\linewidth]{fig6_phase1_novelty.png}
\caption{Objective novelty against the OLMo-2 training corpus: novel 4-grams
per cell, 3 prompts $\times$ 5 seeds. The anti-probable push roughly doubles
surface novelty over min-$p$ sampling.}
\label{fig:phase1}
\end{figure}

\paragraph{Does it climb?} Two tests say no. In verified search (online bin
packing, the FunSearch benchmark, with a deterministic verifier in place of a
judge; evolutionary loops of 5 runs $\times$ 8 generations $\times$ 20--40
candidates on Qwen3-8B and Qwen3-30B-A3B, 3,200 verified heuristics on the
larger model) the anti-probable operator does not separate from plain
sampling: both arms rediscover best-fit at will and neither beats it on held-out
instances. And inside the reverie loop (Section~\ref{sec:loop}), the same
push, against plain sampling in the same scaffold, yields no measurable
difference in judged surprise, connection or novel delta across 1,335
judgments by two judges under three rubrics (the surprise rubric: 2.97 vs
3.48, CI of the difference \ci{-1.50}{+0.49}). Surface novelty is real, and it
does not rise to the level of ideas.

\section{Hypothesis 2: the prompt}
\label{sec:prompt}

We composed three kinds of input (a typical request, a distant concept pair
drawn from a band of embedding distances, and a composed improbable
context) plus two ablations (fragments only, register only), had Claude
Opus~5 develop each into an idea, and had Claude Sonnet~5 judge the result
against its nearest known equivalent ($k=3$; $n\approx15$ per arm). Input
improbability was measured two ways: semantic $k$NN distance to 10,000 real
prompts (all-mpnet-base-v2) and perplexity. The typical prompt matched or beat
every improbable arm, two of which were significantly worse; improbability
did not correlate with judged novelty ($|r|<0.2$, Figure~\ref{fig:prompt});
a pilot effect ($n=8$) had not replicated. Improbable inputs from outside are
noise, at least when a strong model develops them and a strong model judges.

\begin{figure}[t]
\centering
\includegraphics[width=0.72\linewidth]{fig7_improbable_input.png}
\caption{Judged novel delta of the developed idea against the input's
semantic distance from 10,000 real prompts, per arm. No relationship.}
\label{fig:prompt}
\end{figure}
